\phantomsection
\chapter*{Terkenang Dalam Ingatan}
\addcontentsline{toc}{section}{Terkenang Dalam Ingatan --- Aisyah Silviani}
\penulis{Aisyah Silviani}


Kala itu, semua masih terasa indah, semua masih tertata rapi. Setiap hari terasa begitu sederhana, tetapi selalu ada hal kecil yang membuatku merasa bahagia. Sampai suatu kala aku dipanggil untuk pulang. Ternyata ibu memanggilku untuk mandi, karena hari sudah menunjukkan sore hari. Aku pun bergegas pulang dan masuk ke dalam rumah. Di dalam rumah, siapa sangka aku sangat dekat dengan ayah. Ayah yang selalu mengajariku cara untuk menjadi mandiri dan mengajarkanku banyak hal yang mungkin tidak akan pernah aku lupakan.

Dulu, aku sering menghabiskan waktu bersama ayah dan ibu. Kami sering makan bersama, bahkan terkadang hanya dengan satu mangkuk untuk berdua. Hal sederhana seperti itu terasa sangat berarti bagiku. Tidak perlu makanan yang mewah, selama kami bisa duduk bersama dan menikmati makanan sambil bercerita, rasanya sudah cukup membuatku bahagia. Kadang ayah juga mengantarkanku pergi, dan ketika aku pulang, ayah yang menjemputku. Perjalanan yang mungkin terlihat biasa saja, tetapi sekarang menjadi salah satu kenangan yang selalu aku ingat, dan akan selalu kuingat.

Di ruang keluarga, kami juga sering menghabiskan waktu bersama. Aku, ayah, dan ibu duduk bersama sambil bercerita tentang apa saja. Terkadang kami bercanda sampai tertawa bersama. Ada juga saat-saat ketika ayah memberitahuku tentang hal-hal yang tidak boleh aku lakukan. Ayah selalu mengingatkanku agar tidak melakukan sesuatu yang salah dan mengajarkanku untuk mengetahui mana yang boleh dan mana yang tidak boleh dilakukan. Meskipun terkadang terdengar seperti sebuah larangan, sekarang aku mengerti bahwa semua itu adalah bentuk perhatian dan rasa sayang ayah kepadaku.

Mungkin hari-hari itu terlihat sederhana, tetapi ternyata menyimpan begitu banyak kenangan. Makan bersama, pergi diantar, pulang dijemput, bercerita di ruang keluarga, bercanda bersama, sampai nasihat-nasihat kecil dari ayah, semuanya menjadi bagian dari masa yang pernah membuatku merasa sangat lengkap. Kala itu aku tidak pernah berpikir bahwa momen-momen sederhana tersebut akan menjadi sesuatu yang begitu berharga untuk dikenang. Aku hanya menjalani hari seperti biasanya, tanpa menyadari bahwa suatu saat nanti aku akan sangat merindukan semua itu, tak pernah selintas pun di benakku untuk kehilangan setengah dari momen itu. Namun nyatanya, takdir berkata lain.

Malam itu kami masih makan bersama bahkan ayah sempat bercerita tentang hari-hari bagaimana ia bekerja bahkan kami tertawa bersama hanya karena candaan kecil ayah, yang aku pikir momen itu akan bertahan lama namun ternyata itu adalah malam terakhirku bersama ayah. Setelah makan aku memiliki kebiasaan untuk masuk ke kamar dan bermain \textit{handphone}. Saat aku keluar untuk sekadar ke toilet, aku masih melihat ayah masih bermain \textit{handphone} juga di kursi. Aku pikir dia baik-baik saja, namun sepertinya nggak ya, Yah?

Ia pun masuk ke kamar setelah aku selesai dari toilet. Aku mengetahui bahwa ia sudah tidur karena ia harus pergi bekerja di jam 21.30. Aku pikir semua akan berjalan seperti biasa, namun malam itu berbeda. Aku mendengar ibuku memanggil-manggil ayah untuk bangun tapi sepertinya ayah tidak bangun juga. Alhasil ibu memanggilku untuk keluar dari kamar dan aku melihat ayah yang masih terbaring di tempat tidur. Aku mencoba untuk membangunkan ayah dengan segala cara, mulai dari memanggil sampai aku memegang dada dan tangannya untuk membangunkan ayah. ``Ayah, bangun,'' ucapku, namun aku tak kunjung mendapat jawaban.

Tak terasa hatiku mulai terasa berat dan pandanganku mulai buram begitupun dengan suaraku yang mulai bergetar. Aku memikirkan kemungkinan terburuk apa yang terjadi namun aku berusaha untuk menepis pikiranku. Akhirnya aku menelepon kakakku untuk datang ke rumah, tentunya aku menelepon dengan keadaan suara yang bergetar, air mataku sudah mengalir deras. Tak lama kakakku pun datang, ia pun mencoba hal yang sama. Ia berusaha untuk membangunkan ayah dengan segala cara, namun hasilnya nihil ayah tak kunjung bangun juga dan di situ aku mulai menyadari bahwa bibir ayah mulai membiru, detak jantungnya mulai melemah, dan kaki yang mulai dingin.

Tapi aku masih bersyukur karena leher dan tangannya masih hangat saat aku memegangnya dan aku berpikir bahwa masih ada harapan untuk ayah tetap hidup. Pada akhirnya ibuku memanggil \textit{ambulance} dan orang-orang sekitar rumahku membantu untuk menaikkan ayah ke dalam \textit{ambulance}. Aku naik bersama ibu untuk menemani ayah, sepanjang perjalanan aku tidak berhenti menggosok tangan dan kakinya karena anggota tubuhnya mulai dingin. Aku memeluknya dengan keadaan air mataku yang tak kunjung reda. Aku tak berhenti memberikan seluruh kehangatanku untuk ayah yang badannya mulai mendingin.

Tak terasa, kami pun sampai di rumah sakit dan petugas rumah sakit segera membawa ayahku ke IGD. Aku tak bisa berhenti menangis tatkala aku melihat ayah dipasangi segala jenis alat medis di seluruh tubuhnya. Aku pikir masih ada secercah harapan untuk ayah hidup namun$\dots$ dokter menyatakan bahwa ayah telah tiada. Seketika, duniaku terasa runtuh. Aku terjatuh di lantai bersama kakakku dan aku menangis sejadi-jadinya, kemudian aku dibantu untuk bangun kemudian kami memeluk ibu yang masih berusaha tegar di depan kami.

``Gak apa-apa, Allah cuman ngambil apa yang udah jadi milik-Nya,'' ucapnya sambil memeluk kami.

Aku menyadari bahwa yang paling kehilangan adalah ibu, karena ayah adalah \textit{partner} hidup ibu. Namun sebagai anak bungsu dan bisa dibilang aku dekat dengan ayah, sangat wajar kan aku merasa runtuh juga? Hal yang paling menyakitkan adalah ketika ayahku ditutup sekujur tubuhnya di depan mataku sendiri, alat-alat mulai dilepas dari tubuhnya dan kami dianjurkan untuk mengambil surat keterangan kematian.

Saat menunggu, kami sempat berhenti di lorong. Aku membuka kain yang menutupi wajah ayahku, wajahnya terlihat tenang seperti orang yang sedang tertidur. Aku mengusap wajahnya lalu berkata, ``Ayah, nanti aku gimana?'' Aku tak kunjung mendapat jawaban. Aku kembali menangis sejadi-jadinya kemudian aku ditenangkan oleh pamanku. Saat sudah mulai tenang aku melihat ibuku yang masih berusaha tegar di depan anak-anaknya. Aku tak habis pikir, kenapa ibuku bisa melakukannya sedangkan aku seruntuh ini. Aku menghampiri ibuku lalu memeluknya, ia tak berhenti berkata, ``Gak apa-apa masih ada ibu, Allah ngambil apa yang udah jadi titipan.''

Aku tak menjawab apa pun, aku hanya menangis sambil berpikir bagaimana kami kedepannya. Setelah sekian lama menunggu, surat kematian pun telah selesai dibuat dan kami pun pulang. Sepanjang perjalanan, aku tak menyingkirkan tanganku untuk memeluk ayah sambil sesekali air mataku kembali turun. Aku pikir aku sudah selesai menangis, ternyata tidak.

Sesampainya di rumah, tepatnya pada pukul 23.00 rumahku sudah ramai dengan orang-orang. Saat aku mulai melangkah ke dalam rumah, teman dekatku menghampiriku. Aku tak kuasa menahan tangisku, aku jatuh ke lantai sambil memeluknya seraya temanku berkata, ``Yang kuat ya?'' Bahkan ia pun ikut menangis. Aku melihat dengan mata kepalaku sendiri pakaian ayah mulai diganti kain putih dan air mataku tak mau untuk berhenti. Malam itu aku tidur di samping ayah untuk terakhir kalinya, aku memeluk badannya sembari turun air mataku. 

Pada pagi hari, mulai banyak orang-orang yang datang ke rumahku. Mereka menghampiriku dan memelukku seraya berkata, ``Yang kuat ya,'' ``Yang sabar ya.'' Aku tak memberikan jawaban apa pun karena tak kuasa menahan tangis. Tak lama, kami memulai perjalanan untuk mengubur ayahku. Rasanya kakiku lemas tak berdaya saat melihat badannya mulai diturunkan ke dalam tanah. Aku tak bisa menahan tangis tentunya, yang bisa kulakukan hanya mengikhlaskan apa yang sudah Tuhan titipkan. Tak pernah selintas pun di kepalaku, malam itu adalah terakhir kalinya momenku bersama ayah.

Sekarang aku mulai menjalani hari-hariku yang terasa berat. Aku pikir itu hanya momen sesaat namun\(\dots\) sampai sekarang pun ternyata masih berat. Aku kehilangan sosok pembimbingku, yang mengajariku untuk mandiri, yang mengubah pola pikir anak bungsunya yang masih terbilang labil. Aku menjadi anak yang mudah menangis dan sensitif dari sebelumnya, karena sejak hari itu aku belajar mengikhlaskan dengan cara yang terlalu dalam. Kadang kata-kata sederhana bisa menusukku dan rindu selalu datang di waktu yang tidak bisa ditebak. Namun setidaknya aku menyadari bahwa kehilangan tidak membuat seseorang lemah, ia hanya membuat hati lebih rapuh, peka, dan tahu rasanya hancur.

Ayah, maaf jika aku masih berduka atas kehilanganmu.

\(\textit{Based on true story}\)

\jedahias
